% !TeX program = pdflatex
% =============================================================================
% Vorlage: Prüfung / Klausur
% Kantonsschule KFR – Physik
%
% Verwendung:
%   1. Diese Datei nach pruefungen/ kopieren und sinnvoll umbenennen,
%      z. B. pruefungen/kinematik-pruefung.tex
%   2. Kopfzeile (Klasse/Datum/Thema/Titel), Name-Feld, Zeit und Hilfsmittel
%      anpassen.
%   3. Aufgaben mit \question[Punktzahl] ... ergänzen, Lösung jeweils in
%      \begin{solution}...\end{solution} setzen.
%   4. Zweimal kompilieren (pdflatex zweimal laufen lassen), damit die
%      Punktetabelle (\gradetable) und die Bestehensgrenze korrekt
%      berechnet werden.
%   5. Lösungen ein-/ausblenden: \printanswers (Lehrerexemplar) bzw.
%      \noprintanswers (Schülerexemplar, Standard) weiter unten umschalten.
% =============================================================================
\documentclass[addpoints,12pt,a4paper]{exam}

\usepackage[T1]{fontenc}
\usepackage[utf8]{inputenc}
\usepackage[ngerman]{babel}
\usepackage{geometry}
\usepackage{amsmath}
\usepackage{siunitx}
% Hausstil: zusammengesetzte Einheiten immer als echter Bruch (\frac), nicht
% als Inline-Schrägstrich oder \tfrac -- gilt global für jedes \si/\SI.
\sisetup{per-mode=fraction,fraction-command=\frac}
\usepackage{tikz}
\usepackage{physics}
\usepackage{circuitikz}

\geometry{a4paper, margin=2.5cm, headheight=14pt}

% --- Lösungen ein-/ausblenden -----------------------------------------------
% Schülerexemplar (Standard): \noprintanswers
% Lehrer-/Lösungsexemplar:    \printanswers
\noprintanswers
% \printanswers

% --- Punkteformat: "[3 P.]" statt "[3 points]" ------------------------------
\pointname{~P.}
\pointsinmargin
\bonuspointname{~Bonus-P.}

% --- Punktetabelle (\gradetable) auf Deutsch --------------------------------
\hqword{Aufgabe:}
\hpword{Punkte:}
\hsword{Erreicht:}
\htword{Total}

% --- Lösungsbox auf Deutsch --------------------------------------------------
\renewcommand{\solutiontitle}{\noindent\textbf{Lösung:}\enspace}

% --- Kopfzeile: Klasse / Datum / Thema / Titel ------------------------------
% Einheitliches Layout für Arbeitsblatt, Prüfung und Praktikum.
\newcommand{\Kopfzeile}[4]{%
  \begin{center}
    {\Large\bfseries #4}
  \end{center}
  \vspace{0.3em}
  \noindent\textbf{Klasse:}~#1 \hfill \textbf{Datum:}~#2 \hfill \textbf{Thema:}~#3
  \vspace{0.3em}

  \hrule
  \vspace{0.8em}
}

% --- Antwortraum: ca. 2.5 cm Platz pro Punkt für handschriftliche Lösungen --
\newlength{\antwortraumlen}
\newcommand{\Antwortraum}[1]{%
  \pgfmathsetlength{\antwortraumlen}{#1*2.5cm}%
  \vspace{\antwortraumlen}%
  \par
}

% --- Laufender Kopf/Fuss ab Seite 2 -----------------------------------------
\pagestyle{headandfoot}
\firstpageheader{}{}{}
\runningheader{\Thema}{\Klasse\ -- \Datum}{Seite \thepage\ von \numpages}
\runningheadrule
\firstpagefooter{}{}{}
\runningfooter{}{}{}

% Werte für Kopfzeile / laufenden Kopf zentral setzen
\newcommand{\Klasse}{5b}
\newcommand{\Datum}{XX.XX.20XX}
\newcommand{\Thema}{Kinematik}
\newcommand{\PruefungsTitel}{Prüfung: Kinematik}
\newcommand{\Zeit}{45 Minuten}
\newcommand{\Hilfsmittel}{Taschenrechner, Formelsammlung}

\begin{document}

\Kopfzeile{\Klasse}{\Datum}{\Thema}{\PruefungsTitel}

\begin{tabular}{@{}ll@{}}
  \textbf{Name:} \makebox[6cm]{\hrulefill} & \textbf{Zeit:} \Zeit \\[0.6em]
  \textbf{Hilfsmittel:} \Hilfsmittel &
\end{tabular}

\vspace{1em}

% --- Bewertung ---------------------------------------------------------------
% Hinweis: \numpoints wird erst nach zwei Kompilierdurchgängen korrekt
% angezeigt (Standardverhalten der exam-Klasse, siehe Kommentar oben).
\fbox{%
  \parbox{\dimexpr\linewidth-2\fboxsep-2\fboxrule\relax}{%
    \textbf{Bewertung}\\[0.3em]
    Note $= 1 + 3 \times \dfrac{P}{P_{\text{max}}}$\\[0.3em]
    $P$ = erreichte Punkte, \quad $P_{\text{max}}$ = Maximalpunkte\\[0.3em]
    Bestehensgrenze: 60\,\% der Maximalpunkte $=$ Note~4\\[0.5em]
    Maximalpunkte: \numpoints
  }%
}

\vspace{1em}

\gradetable[h][questions]

\vspace{1em}

\begin{questions}

  \question Ein Auto beschleunigt aus dem Stillstand mit
  $a = \SI{2.5}{\meter\per\second\squared}$.

  \begin{parts}
    \part[2] Berechnen Sie die Geschwindigkeit nach $t = \SI{4}{\second}$.
    \Antwortraum{2}

    \part[4] Berechnen Sie die zurückgelegte Strecke nach $t = \SI{4}{\second}$.
    \Antwortraum{4}
  \end{parts}

  \begin{solution}
    \textbf{a)} $v = a \cdot t = \SI{2.5}{\meter\per\second\squared} \cdot \SI{4}{\second} = \SI{10}{\meter\per\second}$

    \textbf{b)} $s = \frac{1}{2} a t^2 = \frac{1}{2} \cdot \SI{2.5}{\meter\per\second\squared} \cdot (\SI{4}{\second})^2 = \SI{20}{\meter}$
  \end{solution}

\end{questions}

\end{document}
